\documentclass[12pt,a4paper]{ctexart}

\usepackage[a4paper,margin=2.5cm]{geometry}
\usepackage{setspace}
\usepackage{array}
\usepackage{longtable}
\usepackage{booktabs}
\usepackage{hyperref}
\usepackage{enumitem}
\usepackage{titlesec}

\hypersetup{
  colorlinks=true,
  linkcolor=black,
  urlcolor=blue,
  citecolor=black
}

\setstretch{1.25}
\setlength{\parindent}{2em}
\titleformat{\section}{\large\bfseries}{\thesection}{1em}{}
\titleformat{\subsection}{\normalsize\bfseries}{\thesubsection}{1em}{}

\title{毕业设计开题报告草稿\\[0.5em]\large 差速驱动移动机器人的运动仿真与控制研究与实现}
\author{学生：刘济伟\\学号：1120220518\\指导教师：郑建波}
\date{\today}

\begin{document}

\maketitle

\section{课题背景与研究意义}
差速驱动移动机器人因结构简单、成本较低、控制实现相对直接，在室内物流、巡检、教学实验和移动平台研究中具有广泛应用。对于此类机器人而言，运动控制性能直接影响路径跟踪精度、姿态稳定性以及整机运行效率，因此围绕其运动仿真、控制算法设计和实物实现开展研究，具有较强的工程意义与应用价值。

在机器人系统开发过程中，仿真平台能够在不增加硬件调试风险的前提下完成模型构建、参数整定和控制策略验证，从而缩短系统开发周期、降低实验成本。对于差速驱动移动机器人而言，仿真环境不仅可以用于验证基本运动学与动力学特性，还能够作为不同控制算法的统一测试平台，为后续实物部署提供依据。

本课题拟围绕差速驱动移动机器人，研究其运动仿真建模方法与控制系统实现方案，构建一套从仿真验证到实物控制的完整技术流程。在控制算法方面，不仅研究传统 PID 控制，还引入 LQR 与 MPC 控制方法，并要求其能够在前期搭建的仿真平台中完成部署与实现。通过对多种控制方法在典型运动任务中的表现进行分析，可为差速驱动移动机器人的控制系统设计提供更系统的实验依据。

\section{国内外研究现状}
差速驱动移动机器人是移动机器人研究中最经典的对象之一，国内外相关研究较为丰富。已有研究主要集中在运动学建模、动力学分析、路径跟踪、轨迹规划和控制算法设计等方面。由于差速驱动结构具有非完整约束特性，其控制问题既具备一定理论研究价值，又便于开展工程实现，因此长期作为机器人控制领域的重要实验平台。

在仿真平台方面，MuJoCo、Gazebo 和 Webots 等工具常被用于移动机器人建模与控制验证。其中，MuJoCo 具有物理求解效率较高、控制接口较灵活、适合快速迭代算法验证等特点，适合用于搭建差速驱动机器人运动仿真平台，并为后续控制算法实验提供统一环境。

在控制算法方面，PID 控制因结构简单、易于工程实现，仍然是移动机器人控制中的常用方法。LQR 控制通过状态空间建模和性能指标优化，能够在系统稳定性与控制代价之间取得较好平衡。MPC 控制则能够显式考虑系统模型与约束条件，在轨迹跟踪和多变量控制场景中表现出较强优势。当前研究趋势是将传统控制方法与模型控制方法结合，在统一平台上进行对比分析，以实现更优的控制效果。

\section{研究目标与研究内容}
\subsection{研究目标}
本课题的总体目标是围绕差速驱动移动机器人，完成运动仿真平台搭建、控制算法设计与实物系统实现，具体目标如下：
\begin{enumerate}[label=\arabic*.]
  \item 完成差速驱动移动机器人的运动仿真模型构建与仿真环境搭建；
  \item 设计并实现 PID、LQR 和 MPC 三类控制算法，并在仿真平台中完成部署与验证；
  \item 实现仿真平台与实物机器人控制链路的基本对接，完成控制系统实现；
  \item 对不同控制算法在典型运动任务中的控制效果进行对比分析；
  \item 总结差速驱动移动机器人仿真与控制实现过程中的关键问题与改进方向。
\end{enumerate}

\subsection{研究内容}
围绕上述目标，本文拟开展以下几个方面的工作：
\begin{enumerate}[label=\arabic*.]
  \item 差速驱动移动机器人的运动学与动力学分析；
  \item 基于 MuJoCo 的差速驱动移动机器人仿真平台搭建；
  \item PID 控制算法设计与仿真实现；
  \item LQR 控制算法建模、求解与仿真实现；
  \item MPC 控制算法建模、约束设计与仿真实现；
  \item 实物机器人控制系统实现与仿真结果对比分析。
\end{enumerate}

\section{技术路线与实现方案}
本课题采用“先建模仿真、后算法部署、再实物验证”的技术路线。首先对差速驱动移动机器人进行运动学分析，明确线速度、角速度与左右轮速度之间的关系，并据此建立仿真模型。随后在 MuJoCo 平台中完成机器人模型、场景和控制接口设计，形成可重复测试的仿真实验环境。

在控制算法实现阶段，首先完成 PID 控制器设计，用于实现基础速度控制与姿态调节。在此基础上，进一步建立状态空间模型，完成 LQR 控制器设计与参数整定，并分析其在稳定性和控制能耗方面的特点。随后构建 MPC 控制器，通过预测模型和滚动优化机制实现更高质量的轨迹跟踪控制。三类控制算法均要求在前期仿真平台中完成实际部署和运行，而不是仅停留在理论推导层面。

在系统实现阶段，将仿真验证过的控制逻辑迁移到实物机器人控制程序中，完成控制指令生成、底层驱动与基本运动验证。最终以直线运动、转向控制和轨迹跟踪等任务为测试场景，对不同控制算法的响应速度、稳态误差和控制平滑性进行比较分析。

\section{研究基础与可行性分析}
课题已有一定前期基础。此前已完成“基于 Linux 的简易差速驱动机器人运动仿真实现”相关研究，已经在差速驱动机器人平台上完成基础仿真与控制实验，对 Linux 开发环境、MuJoCo 仿真平台、差速驱动运动控制以及控制参数整定积累了一定经验。这些前期工作为本课题继续深入开展差速驱动移动机器人的仿真与控制研究提供了基础。

在硬件条件方面，当前具备可用于实验验证的差速驱动移动机器人实物平台，以及 Jetson Orin Nano 等计算资源，可支持控制算法部署与调试。在软件条件方面，MuJoCo、Python 及相关控制算法开发工具链较为成熟，能够满足本课题的实现需求。

从研究内容上看，PID 控制具备较成熟的实现基础，LQR 与 MPC 虽然在建模和调参方面要求更高，但由于研究对象为差速驱动移动机器人，其系统结构相对清晰、状态维度可控，适合在本科毕业设计阶段开展建模、实现与对比分析。因此，本课题总体具有较好的可行性。

\section{进度安排}
根据当前毕业设计时间节点，计划安排如下：

\begin{longtable}{p{3cm}p{4cm}p{7cm}}
\toprule
阶段 & 时间安排 & 主要内容 \\
\midrule
开题准备 & 2026-03-25 至 2026-03-27 & 明确课题方案，完成开题报告与总体技术路线设计 \\
仿真平台搭建 & 2026-03-28 至 2026-04-05 & 完成差速驱动机器人模型构建与 MuJoCo 仿真环境搭建 \\
控制算法实现 & 2026-04-06 至 2026-04-20 & 完成 PID、LQR、MPC 控制算法的建模、实现与仿真部署 \\
中期检查 & 2026-04 下旬 & 展示仿真平台、控制算法效果及阶段性分析结果 \\
系统完善与实验 & 2026-05-01 至 2026-05-15 & 完成实物控制实现及多组实验对比分析 \\
论文撰写与答辩准备 & 2026-05 中下旬至 2026-06 上旬 & 完成论文撰写、修改、答辩材料整理与演示准备 \\
\bottomrule
\end{longtable}

\section{预期成果}
本课题预期形成以下成果：
\begin{enumerate}[label=\arabic*.]
  \item 差速驱动移动机器人的运动仿真平台；
  \item PID、LQR、MPC 三类控制算法的实现程序与仿真验证结果；
  \item 差速驱动移动机器人实物控制系统的基础实现；
  \item 不同控制算法的实验对比分析结果；
  \item 一篇完整的毕业设计论文。
\end{enumerate}

\section{可能存在的问题与解决思路}
在研究实施过程中，可能面临以下问题：
\begin{enumerate}[label=\arabic*.]
  \item 仿真模型与实物参数之间存在偏差，影响算法迁移效果；
  \item LQR 与 MPC 的建模和参数整定工作量较大；
  \item MPC 控制计算量相对较高，可能影响实时性；
  \item 毕设时间较紧，多种控制算法同时实现存在进度压力。
\end{enumerate}

针对上述问题，可采用“先完成 PID 基础闭环控制，再逐步扩展到 LQR 和 MPC”“先在仿真平台中完成统一测试，再选择适合的方案部署实物”“通过简化模型与合理约束降低 MPC 计算复杂度”“以阶段性里程碑推进任务”等方式，保证课题主体内容按期完成。

\end{document}
